\documentclass[conference]{IEEEtran}

\usepackage[T1]{fontenc}
\usepackage{amsmath,amssymb}
\usepackage{graphicx}
\usepackage{booktabs}
\usepackage{listings}
\usepackage[hidelinks]{hyperref}

\title{A Clear and Specific Paper Title}

\author{
  \IEEEauthorblockN{First Author}
  \IEEEauthorblockA{Department or Laboratory \\
  Institution \\
  City, Country \\
  author@example.edu}
}

\begin{document}

\maketitle

\begin{abstract}
State the problem, method, principal result, and implication in a compact
paragraph.
\end{abstract}

\begin{IEEEkeywords}
scientific writing, reproducibility, research software
\end{IEEEkeywords}

\section{Introduction}
\label{sec:introduction}

Explain the research question and position it against relevant work such as
\cite{lamport1994latex}.

\section{Method}
\label{sec:method}

Define the method, assumptions, and evaluation protocol.

\begin{equation}
  \label{eq:objective}
  \mathcal{L}(\theta) = \frac{1}{n}\sum_{i=1}^{n}
  \ell\bigl(f_{\theta}(x_i), y_i\bigr).
\end{equation}

\section{Results}
\label{sec:results}

\begin{table}[t]
  \centering
  \caption{Replace this table with the primary comparison.}
  \label{tab:results}
  \begin{tabular}{lrr}
    \toprule
    Method & Accuracy & Time (s) \\
    \midrule
    Baseline & 0.81 & 14.2 \\
    Proposed & 0.86 & 12.8 \\
    \bottomrule
  \end{tabular}
\end{table}

Report uncertainty and all material experimental conditions.

\section{Conclusion}

Summarize the supported claim and limitations.

\bibliographystyle{IEEEtran}
\bibliography{references}

\end{document}
